% =====================================================================
% config/styles.tex
% Warna, tata letak, formato judul, header/footer, dan penomoran.
% =====================================================================

% --- Palet warna ------------------------------------------------------
\definecolor{ubbblue}{RGB}{15,55,110}
\definecolor{ubbteal}{RGB}{0,120,130}
\definecolor{ubbgold}{RGB}{200,150,40}
\definecolor{codebg}{RGB}{247,247,247}
\definecolor{codeframe}{RGB}{210,215,225}
\definecolor{textgray}{RGB}{90,90,90}

% --- Spasi dan paragraf ------------------------------------------------
\setlength{\parindent}{0pt}
\setlength{\parskip}{0.5em}
\linespread{1.12}
\setlength{\emergencystretch}{2em}
\raggedbottom
\clubpenalty=1000
\widowpenalty=1000
\renewcommand{\arraystretch}{1.15}

% --- Format judul bab --------------------------------------------------
\titleformat{\chapter}[display]
  {\normalfont\sffamily\bfseries\color{ubbblue}}
  {\Large\MakeUppercase{\chaptertitlename}~\thechapter}
  {10pt}
  {\Huge}
\titlespacing*{\chapter}{0pt}{-10pt}{28pt}

% --- Format judul section ---------------------------------------------
\titleformat{\section}
  {\normalfont\Large\bfseries\color{ubbblue}}
  {\thesection}{0.6em}{}
\titlespacing*{\section}{0pt}{2.0ex plus 0.5ex minus 0.2ex}{1.0ex}

\titleformat{\subsection}
  {\normalfont\large\bfseries\color{ubbteal}}
  {\thesubsection}{0.5em}{}
\titlespacing*{\subsection}{0pt}{1.5ex plus 0.3ex minus 0.2ex}{0.8ex}

% --- Daftar isi --------------------------------------------------------
\setcounter{tocdepth}{2}
\setlength{\cftchapnumwidth}{2.8em}
\renewcommand{\cftchapfont}{\bfseries}
\renewcommand{\cftchappagefont}{\bfseries}
\renewcommand{\cftchapleader}{\cftdotfill{\cftdotsep}}

% --- Keterangan gambar/tabel -----------------------------------------
\captionsetup{
  font=small,
  labelfont={bf,color=ubbblue},
  labelsep=period,
  skip=6pt
}

% --- Header dan footer -------------------------------------------------
\pagestyle{fancy}
\fancyhf{}
\fancyhead[L]{\small\itshape\color{textgray}Teknik Kompilasi}
\fancyhead[R]{\small\color{textgray}\thepage}
\renewcommand{\headrulewidth}{0.4pt}
\renewcommand{\footrulewidth}{0pt}

\fancypagestyle{plain}{%
  \fancyhf{}%
  \fancyfoot[C]{\small\thepage}%
  \renewcommand{\headrulewidth}{0pt}%
}

% --- Warna tautan ------------------------------------------------------
\hypersetup{
  colorlinks=true,
  linkcolor=ubbblue,
  citecolor=ubbteal,
  urlcolor=ubbteal,
  filecolor=ubbteal
}

% --- Penomoran ---------------------------------------------------------
% Bab memakai angka Romawi (Bab I, Bab II, ...).
% Section, gambar, tabel, kode, dan rumus memakai nomor bab Arab.
\renewcommand{\thechapter}{\Roman{chapter}}
\renewcommand{\thesection}{\arabic{chapter}.\arabic{section}}
\renewcommand{\thesubsection}{\thesection.\arabic{subsection}}
\renewcommand{\thefigure}{\arabic{chapter}.\arabic{figure}}
\renewcommand{\thetable}{\arabic{chapter}.\arabic{table}}
\renewcommand{\theequation}{\arabic{chapter}.\arabic{equation}}

% --- Paragraf yang lebih rapi di kolom sempit -------------------------
\setlength{\skip\footins}{10pt}
